\chapter{Neutropenia (Low Neutrophil Count)}

\section*{Definition}
Neutropenia is a reduced absolute neutrophil count (ANC), the blood-cell measure most closely related to bacterial and fungal infection risk. Risk generally increases as the ANC falls, especially below 500 cells per microliter or when the decline is rapid.

\section*{When to See a Doctor}
A temperature of 38 degrees C (100.4 degrees F) or higher in a person with severe or treatment-related neutropenia is an emergency. Seek immediate care for fever, chills, new cough, shortness of breath, painful urination, abdominal pain, mouth sores, confusion, or rapidly worsening illness.

\section*{How It Develops}
Neutrophils provide early defense against bacterial and fungal invasion. A low count reduces the inflammatory response, so serious infection may cause little local redness or pus while progressing rapidly.

\section*{Causes}
\begin{itemize}
  \item Chemotherapy, radiation, marrow infiltration, myelodysplastic disease, leukemia, and severe infection.
  \item Medicines including antithyroid drugs, clozapine, some antibiotics, anticonvulsants, and immunosuppressants.
  \item Autoimmune neutropenia, nutritional deficiency, congenital syndromes, and hypersplenism.
  \item Transient viral suppression, including influenza, EBV, CMV, and other viral illnesses.
\end{itemize}

\section*{Related Symptoms}
\begin{itemize}
  \item Neutropenia itself may cause no symptoms; fever or infection is often the first sign.
  \item Recurrent mouth ulcers, skin infections, pneumonia, sinusitis, or perianal pain may occur.
\end{itemize}

\section*{Medical Evaluation}
\begin{itemize}
  \item Repeat complete blood count with differential and review the trend, prior counts, medicines, chemotherapy, infections, and family history.
  \item Examine for infection and obtain cultures, imaging, and targeted viral or autoimmune testing when indicated.
  \item Persistent, severe, unexplained, or multilineage cytopenia may require hematology review and bone-marrow examination.
\end{itemize}

\section*{Treatment Options}
Febrile neutropenia requires prompt broad-spectrum antibiotics after cultures when this does not delay treatment. Management may include stopping a causative drug, granulocyte colony-stimulating factor in selected patients, antifungal therapy for persistent high-risk fever, and treatment of the underlying marrow or immune disorder.

\section*{Self-Care and Home Management}
Follow the oncology or hematology action plan. Check temperature when unwell, practice careful hand hygiene, avoid sick contacts and unsafe food, and do not delay emergency assessment for fever. Avoid rectal thermometers or procedures unless specifically advised.

\section*{References}
\begin{thebibliography}{99}
\bibitem{neutropenia:ref1} Freifeld AG, et al. Clinical practice guideline for the use of antimicrobial agents in neutropenic patients with cancer. \textit{Clin Infect Dis}. 2011;52:e56--e93.
\bibitem{neutropenia:ref2} National Comprehensive Cancer Network. Prevention and treatment of cancer-related infections. NCCN Guidelines; accessed 2025.
\end{thebibliography}
